\section{Introduction and current theorem structure}\label{sec:current-overview}\label{sec:introduction}
A finite-dimensional prepared quantum system can have an optimal observable probability response whose exact classical implementation needs arbitrarily large memory. The required size depends not only on noise magnitude but on rational approximation of its phase. We distinguish the whole memory curve, its logarithmic growth exponents and a uniform power law. These are different mathematical questions even at one fixed phase.

The new main theorem, Theorem~\ref{thm:r71-envelope}, identifies an explicit arithmetic envelope comparable to the least exact stochastic memory with universal constants. It has no badly-approximable hypothesis. Theorem~\ref{thm:r71-exponents} determines both logarithmic exponents from the irrationality exponent of any fixed phase. Theorem~\ref{thm:r71-cubic-rigidity} proves that a uniform cubic-root law holds exactly for badly approximable phases, and proves unbounded fluctuations after cubic-root normalization at every other irrational phase. Corollary~\ref{cor:r71-mixtures} carries the envelope to finite spectral mixtures and leakage. All new statements have full proofs below.

The response-to-spectrum equivalence, stochastic sector inequality and its equality structure, rational polygon construction and exact leakage increment are inherited from the preceding complete manuscript. They are included after the new theorem to expose its complete positive-matrix foundation. Classical stochastic eigenvalue regions and positive-realization methods are not being renamed as new inventions \cite{DmitrievDynkin46,JohnsonPaparella17,KirklandSmigoc22,MonrasWinter14}. The continued-fraction identity used here is classical and is proved explicitly alongside its credit \cite{HoneVarona19}. Noise-dependent lower bounds for classical realizations of quantum processes, including a worst-angle square-root law, precede this paper \cite{Fanizza24}; the present conclusions fix an arbitrary phase and determine its entire approximation-sensitive envelope and two exact growth exponents.

The operational interpretation is unchanged. A real three-outcome audit with one fixed squared loss and two fixed decoder values forces the target probability response at the original least expected-cost boundary. No extra observation of the protected system, free initial mixture, uncharged clock or freely retained random seed is used. The new theorem does not infer a classical exact-memory rate merely from the physical Hilbert dimension or physical mixing speed.

The full current supplement preserves the stopping-rigidity, robust-tie, instrument-realization, finite-model witness and performed-law theorems and the original reset, rank, spectral, geometric, predictive and effective developments with their individual assumptions. The main article also retains the actual-history lower and common-program upper proof. The source retention map separates this paired current mathematics from all unchanged historical manuscripts and independent realization lines. A new phase theorem neither silently deletes those results nor makes every independent companion a premise of its proof.
